\scanpage{594}

\mspage{3}

\unsure{-} Soit

\[
\unsure{f_{A,\Phi}|\varphi_{A,\Phi}}\in k_\Phi
\]

l'image de $f_A$ par \inserted{l'hom. can.} $k_E\longrightarrow k_\Phi$. Si $E=A\amalg B$ on aura

\[
f_A+f_B=1,\quad f_A f_B=0
\]

d'où

\[
f_{A,\Phi}+f_{B,\Phi}=1,\quad f_{A,\Phi}f_{B,\Phi}=0\ .
\]

Si $k_\Phi$ est un corps, cela implique que \deleted{$\unsure{f_A J}$} on a $f_{A,\Phi}=0$ ou $f_{B,\Phi}=0$. Mais si p.ex. $f_{A,\Phi}=0$, on voit que cel\unsure{a} signifie que il existe un $J\in\Phi$ tel que $f_A\mid J=0$, ce qui signifie aussi que $J\subset B$. De $=$, si $f_{B,\Phi}=0$, on aura $J\subset A$. Donc $\Phi$ est un ultrafiltre.

Inversement, supposons que $\Phi$ soit un ultrafiltre. \deleted{Alors} Prouvons que $k_\Phi$ est un corps, i.e. que pour $f\in k_\Phi$, $f\ne0\Rightarrow f$ inversible.

\deleted{Mais $f=0$ signifie qu'il} Notons que $f$ provient d'un $\unsure{g}$ \unsure{élmt|él'mt} d'un $k_{\deleted{\unsure{B|D}}J}$ $(J\in\Phi)$, donc est l'image canonique d'un $g\in k_E$. Dire que $f=0$ signifie que $\exists\ J\in\Phi$ tel que $g\mid J=0$. L'inverse \inserted{$f\ne0$} signifie donc que pour tt $J\in\Phi$, on a $g\mid J\ne0$.

% archive sheet number: 593
